\documentclass[11pt,a4paper]{article}
\usepackage[utf8]{inputenc}
\usepackage[T1]{fontenc}
\usepackage{lmodern}
\usepackage[margin=2.5cm]{geometry}
\usepackage{hyperref}

\title{What Does a State-Space Model Contribute to Event-Camera Object Detection?\\
\large A Controlled Study with a Spiking Variant}
\author{Benas Vaiciulis\\UNSW Sydney}
\date{}

\begin{document}

\maketitle

\begin{abstract}
\end{abstract}

\section{Introduction}
\subsection{Event Cameras and the Detection Problem}
\subsection{Research Question}
\subsection{Contributions}

\section{Related Work}
\subsection{Event-Based Object Detection}
\subsection{State-Space Models and Mamba}
\subsection{State-Space Models for Event Cameras}
\subsection{Spiking Neural Networks for Event-Based Detection}

\section{Method}
\subsection{Event Representation}
\subsection{The Frozen Detection Pipeline}
\subsection{Spatial and Temporal Mixing Axes}
\subsection{EventSSM: Convolutional Spatial, State-Space Temporal}
\subsection{PureSSM: State-Space Spatial and Temporal}
\subsection{SpikingSSM: A Spiking Temporal Readout}
\subsection{The Readout Ladder: Dynamics, Sparsity and Binarisation}

\section{Experimental Setup}
\subsection{Dataset and Metrics}
\subsection{Training Protocol}
\subsection{Baseline Reproduction}
\subsection{Evaluation Protocol}
\subsection{Hardware and Implementation}

\section{Results}
\subsection{Accuracy}
\subsection{Effective Receptive Field}
\subsection{Event-Rate Robustness}
\subsection{Reproducing Inference-Time \texorpdfstring{$\Delta t$}{Delta t}-Rescaling}
\subsection{Efficiency}
\subsection{The Cost of Spiking}

\section{Discussion}
\subsection{What the State-Space Model Contributes}
\subsection{Latency and the Choice of Window}
\subsection{The Neuromorphic Case}
\subsection{Scope Change from UAV Integration}
\subsection{Limitations}

\section{Conclusion and Future Work}
\subsection{Conclusion}
\subsection{Future Work}

\section*{Acknowledgments}

\section*{References}

\end{document}
